\chapter{Validation et résultats}
\section{Objectif de l'évaluation}
L'évaluation a été conçue pour distinguer deux questions : le système a-t-il proposé le bon champ, et la valeur associée est-elle correcte ? Cette distinction évite d'interpréter un taux de remplissage élevé comme une exactitude automatique. L'évaluateur découvre les quadruplets composés d'un rapport, d'un gabarit vide, d'un gabarit renseigné et d'une prédiction. Il vérifie la structure JSON, aplatit les chemins complets, normalise les unités et les textes, puis calcule des agrégats micro et macro.

Les indicateurs produits sont la précision, le rappel, le F1, le taux d'autoremplissage, le faux autoremplissage, le taux de validation et l'exactitude des valeurs. Pour la recherche, un rappel@k indicatif peut également être calculé. L'évaluation est exécutée hors ligne et peut imposer des seuils de qualité dans l'intégration continue.

\section{Jeux de données et précautions d'interprétation}
Les résultats présentés proviennent de jeux synthétiques ou de données localement préparées pour l'évaluation. Ils mesurent la capacité de la chaîne sur les formats testés. Ils ne démontrent ni une performance sur un hôpital indépendant, ni une aptitude au diagnostic, ni une validation clinique. Cette limitation est centrale : une généralisation nécessiterait des jeux représentatifs, des annotations contrôlées, une analyse de biais et une procédure de validation conforme au contexte d'usage.

\section{Résultats sur les bilans biologiques}
Sur un sous-ensemble comparable de 120 échantillons, les améliorations du parsing et de la reconstruction des lignes ont augmenté le rappel de 0,8086 à 0,9569. Le F1 est passé de 0,8942 à 0,9780, sans faux autoremplissage observé. La précision est restée à 1,0000. L'exactitude des valeurs a légèrement diminué de 0,9452 à 0,9320 : la récupération de valeurs auparavant manquées expose aussi davantage de divergences de valeur à comparer.

\begin{table}[H]
\centering
\caption{Résultats biologiques avant et après les améliorations de rappel (120 échantillons comparables)}
\begin{tabular}{lcc}
\toprule
Indicateur & Avant & Après\\
\midrule
Précision & 1,0000 & 1,0000\\
Rappel & 0,8086 & 0,9569\\
F1 & 0,8942 & 0,9780\\
Exactitude des valeurs & 0,9452 & 0,9320\\
Faux autoremplissage & 0,0000 & 0,0000\\
\bottomrule
\end{tabular}
\end{table}

\begin{figure}[H]
\centering
\begin{tikzpicture}
\begin{axis}[ybar,bar width=16pt,width=.82\textwidth,height=7cm,ymin=0,ymax=1.08,ylabel={Score},symbolic x coords={Rappel,F1,Exactitude},xtick=data,legend style={at={(0.5,-0.18)},anchor=north,legend columns=2},nodes near coords, every node near coord/.append style={font=\scriptsize}]
\addplot[fill=ensiblue] coordinates {(Rappel,0.8086) (F1,0.8942) (Exactitude,0.9452)};
\addplot[fill=green!55!black] coordinates {(Rappel,0.9569) (F1,0.9780) (Exactitude,0.9320)};
\legend{Avant,Après}
\end{axis}
\end{tikzpicture}
\caption{Évolution des résultats biologiques}
\end{figure}

Sur l'évaluation finale de 1 000 échantillons biologiques, les valeurs reportées sont une précision de 1,0000, un rappel de 0,9425, un F1 de 0,9704 et une exactitude des valeurs de 0,9264. Ces chiffres confirment le gain de couverture tout en indiquant qu'un contrôle des valeurs reste nécessaire.

\section{Résultats sur les comptes rendus d'imagerie}
Un premier jeu synthétique simple a fourni une exactitude des valeurs de 1,0000 et un F1 de 0,9962. Ce résultat est volontairement considéré comme optimiste, car les rapports et les gabarits partageaient une formulation régulière. Un second jeu plus difficile a introduit des variations cm/mm, des mesures de base et courantes, des lésions non cibles, des distracteurs et des conclusions RECIST potentiellement contradictoires avec des mentions précédentes.

\begin{table}[H]
\centering
\caption{Résultats du jeu d'imagerie difficile sur 152 échantillons terminés}
\begin{tabular}{lcc}
\toprule
Indicateur & Avant corrections ciblées & Après corrections ciblées\\
\midrule
Précision & 1,0000 & 1,0000\\
Rappel & 0,9671 & 1,0000\\
F1 & 0,9833 & 1,0000\\
Exactitude des valeurs & 0,6395 & 0,9901\\
Faux autoremplissage & 0,0000 & 0,0000\\
\bottomrule
\end{tabular}
\end{table}

Les corrections ont porté sur la normalisation d'unités, la sélection de la mesure de base et la priorité donnée à la conclusion RECIST. Après correction, la taille de base était correcte dans 152 cas sur 152; la réponse RECIST l'était dans 149 cas sur 152. Les trois erreurs restantes correspondent à deux confusions \emph{CR vers SD} et une confusion \emph{CR vers PR}. L'extraction planifiée pour 200 rapports a été interrompue après 152 traitements complets; les métriques portent donc exclusivement sur ces 152 quadruplets complets.

\section{Analyse des résultats}
Les résultats montrent que la principale amélioration ne vient pas d'un unique modèle plus complexe, mais d'une meilleure conservation du contexte et de règles ciblées. Pour les bilans, la reconstruction de lignes complètes, la prise en charge de formats supplémentaires et le mapping d'alias ont amélioré le rappel. Pour l'imagerie, conserver le rapport court entier et hiérarchiser les indices de temporalité ont amélioré l'exactitude des valeurs.

Le F1 et l'exactitude des valeurs doivent être lus ensemble. L'état initial de l'imagerie difficile constitue un exemple instructif : le F1 était déjà de 0,9833 alors que l'exactitude des valeurs n'était que de 0,6395. Le système remplissait donc souvent les champs attendus, mais ne sélectionnait pas toujours la bonne mesure ou la bonne conclusion.

\section{Limites et perspectives}
Plusieurs limites restent identifiées. Le projet ne remplace pas une revue humaine avant écriture définitive. Les jeux de données utilisés ne sont pas une validation externe. L'OCR de documents scannés, la pseudonymisation avant indexation, une API HTTP sécurisée, la journalisation opérationnelle, le suivi de latence et le monitoring de dérive constituent des évolutions pertinentes. Une étude future devrait aussi évaluer la robustesse à une plus grande variété d'établissements, de langues et de mises en page.

\section{Conclusion}
La validation confirme l'intérêt d'une approche spécialisée, traçable et testable. Les gains observés sont utiles pour guider le développement, mais leur portée est explicitement limitée au cadre expérimental décrit.
